% Generated by experiments/make_numbers.py -- do not edit.
%
% Rejection rates of every matched family read in each of the three
% tails, from results/grid_A_B_tail_robustness.csv and
% results/grid_C_tail_robustness_by_effect.csv.

\begin{table}[t]
\centering
\small
\begin{tabular}{llrrr}
\toprule
Effect & Matched family & Lower & Two-sided & Upper \\
\midrule
\multicolumn{5}{l}{\emph{Grids A and B, no planted effect ($n = 1,680$ runs)} --- these are false-positive rates} \\
 & random-matched & 0.035 & 0.199 & 0.210 \\
 & expression-matched & 0.060 & 0.067 & 0.058 \\
 & codetection-matched & 0.079 & 0.076 & 0.043 \\
\midrule
\multicolumn{5}{l}{\emph{Grid C, planted effect ($n = 240$ runs per effect)} --- these are power} \\
$0.25$ & random-matched & 0.000 & 0.779 & 0.808 \\
 & expression-matched & 0.000 & 0.871 & 0.908 \\
 & codetection-matched & 0.000 & 0.863 & 0.900 \\
$0.5$ & random-matched & 0.000 & 0.900 & 0.900 \\
 & expression-matched & 0.000 & 0.988 & 0.992 \\
 & codetection-matched & 0.000 & 0.988 & 0.996 \\
\bottomrule
\end{tabular}
\caption{The matched null in three conventions, computed from the same
null draws in every case. The paper reports the two-sided column; the
directional columns are given so that the choice can be checked rather
than taken on trust. Reading them together is what separates the two
ways a matched null can fail: a false-positive rate that depends on the
tail is a defect of the test, while a power that the lower tail cannot
reach is not a defect at all but the reason the convention has to be
stated before the grid is run.}
\label{tab:tail}
\end{table}
